% Options for packages loaded elsewhere
% Options for packages loaded elsewhere
\PassOptionsToPackage{unicode}{hyperref}
\PassOptionsToPackage{hyphens}{url}
\PassOptionsToPackage{dvipsnames,svgnames,x11names}{xcolor}
%
\documentclass[
  12pt,
]{article}
\usepackage{xcolor}
\usepackage[margin=2.5cm,bottom=3cm]{geometry}
\usepackage{amsmath,amssymb}
\usepackage{cancel}
\setcounter{secnumdepth}{5}
\usepackage{iftex}
\ifPDFTeX
  \usepackage[T1]{fontenc}
  \usepackage[utf8]{inputenc}
  \usepackage{textcomp} % provide euro and other symbols
\else % if luatex or xetex
  \usepackage{unicode-math} % this also loads fontspec
  \defaultfontfeatures{Scale=MatchLowercase}
  \defaultfontfeatures[\rmfamily]{Ligatures=TeX,Scale=1}
\fi
\usepackage{lmodern}
\ifPDFTeX\else
  % xetex/luatex font selection
\fi
% Use upquote if available, for straight quotes in verbatim environments
\IfFileExists{upquote.sty}{\usepackage{upquote}}{}
\IfFileExists{microtype.sty}{% use microtype if available
  \usepackage[]{microtype}
  \UseMicrotypeSet[protrusion]{basicmath} % disable protrusion for tt fonts
}{}
\makeatletter
\@ifundefined{KOMAClassName}{% if non-KOMA class
  \IfFileExists{parskip.sty}{%
    \usepackage{parskip}
  }{% else
    \setlength{\parindent}{0pt}
    \setlength{\parskip}{6pt plus 2pt minus 1pt}}
}{% if KOMA class
  \KOMAoptions{parskip=half}}
\makeatother
% Make \paragraph and \subparagraph free-standing
\makeatletter
\ifx\paragraph\undefined\else
  \let\oldparagraph\paragraph
  \renewcommand{\paragraph}{
    \@ifstar
      \xxxParagraphStar
      \xxxParagraphNoStar
  }
  \newcommand{\xxxParagraphStar}[1]{\oldparagraph*{#1}\mbox{}}
  \newcommand{\xxxParagraphNoStar}[1]{\oldparagraph{#1}\mbox{}}
\fi
\ifx\subparagraph\undefined\else
  \let\oldsubparagraph\subparagraph
  \renewcommand{\subparagraph}{
    \@ifstar
      \xxxSubParagraphStar
      \xxxSubParagraphNoStar
  }
  \newcommand{\xxxSubParagraphStar}[1]{\oldsubparagraph*{#1}\mbox{}}
  \newcommand{\xxxSubParagraphNoStar}[1]{\oldsubparagraph{#1}\mbox{}}
\fi
\makeatother


\usepackage{longtable,booktabs,array}
\newcounter{none} % for unnumbered tables
\usepackage{calc} % for calculating minipage widths
% Correct order of tables after \paragraph or \subparagraph
\usepackage{etoolbox}
\makeatletter
\patchcmd\longtable{\par}{\if@noskipsec\mbox{}\fi\par}{}{}
\makeatother
% Allow footnotes in longtable head/foot
\IfFileExists{footnotehyper.sty}{\usepackage{footnotehyper}}{\usepackage{footnote}}
\makesavenoteenv{longtable}
\usepackage{graphicx}
\makeatletter
\newsavebox\pandoc@box
\newcommand*\pandocbounded[1]{% scales image to fit in text height/width
  \sbox\pandoc@box{#1}%
  \Gscale@div\@tempa{\textheight}{\dimexpr\ht\pandoc@box+\dp\pandoc@box\relax}%
  \Gscale@div\@tempb{\linewidth}{\wd\pandoc@box}%
  \ifdim\@tempb\p@<\@tempa\p@\let\@tempa\@tempb\fi% select the smaller of both
  \ifdim\@tempa\p@<\p@\scalebox{\@tempa}{\usebox\pandoc@box}%
  \else\usebox{\pandoc@box}%
  \fi%
}
% Set default figure placement to htbp
\def\fps@figure{htbp}
\makeatother





\setlength{\emergencystretch}{3em} % prevent overfull lines

\providecommand{\tightlist}{%
  \setlength{\itemsep}{0pt}\setlength{\parskip}{0pt}}



 
\usepackage[style=apa,]{biblatex}
\addbibresource{references.bib}


% Page look
\usepackage{fancyhdr}
\pagestyle{fancy}
\fancyhf{} 
\renewcommand{\headrulewidth}{0pt} % Removes the horizontal line at the top
% Page numbering
\fancyfoot[R]{\thepage}           % Puts page number in the Right footer
% This forces the "plain" style (used by first pages) 
% to also put the page number on the right
\fancypagestyle{plain}{
    \fancyhf{}
    \fancyfoot[R]{\thepage}
    \renewcommand{\headrulewidth}{0pt}
}

% Math macros
\newcommand{\indep}{\perp \!\!\! \perp}

\usepackage{cancel}
\usepackage{algorithm}
\usepackage{algpseudocode}
\usepackage{lscape}
\makeatletter
\@ifpackageloaded{caption}{}{\usepackage{caption}}
\AtBeginDocument{%
\ifdefined\contentsname
  \renewcommand*\contentsname{Table of contents}
\else
  \newcommand\contentsname{Table of contents}
\fi
\ifdefined\listfigurename
  \renewcommand*\listfigurename{List of Figures}
\else
  \newcommand\listfigurename{List of Figures}
\fi
\ifdefined\listtablename
  \renewcommand*\listtablename{List of Tables}
\else
  \newcommand\listtablename{List of Tables}
\fi
\ifdefined\figurename
  \renewcommand*\figurename{Figure}
\else
  \newcommand\figurename{Figure}
\fi
\ifdefined\tablename
  \renewcommand*\tablename{Table}
\else
  \newcommand\tablename{Table}
\fi
}
\@ifpackageloaded{float}{}{\usepackage{float}}
\floatstyle{ruled}
\@ifundefined{c@chapter}{\newfloat{codelisting}{h}{lop}}{\newfloat{codelisting}{h}{lop}[chapter]}
\floatname{codelisting}{Listing}
\newcommand*\listoflistings{\listof{codelisting}{List of Listings}}
\makeatother
\makeatletter
\usepackage{pdflscape}
\makeatother
\makeatletter
\makeatother
\makeatletter
\@ifpackageloaded{caption}{}{\usepackage{caption}}
\@ifpackageloaded{subcaption}{}{\usepackage{subcaption}}
\makeatother
\usepackage{bookmark}
\IfFileExists{xurl.sty}{\usepackage{xurl}}{} % add URL line breaks if available
\urlstyle{same}
\hypersetup{
  pdftitle={Quantifying the Unknown},
  colorlinks=true,
  linkcolor={blue},
  filecolor={Maroon},
  citecolor={Blue},
  urlcolor={Blue},
  pdfcreator={LaTeX via pandoc}}


\title{Quantifying the Unknown}
\usepackage{etoolbox}
\makeatletter
\providecommand{\subtitle}[1]{% add subtitle to \maketitle
  \apptocmd{\@title}{\par {\large #1 \par}}{}{}
}
\makeatother
\subtitle{A numerical analysis of surprise-based observability of
acceptance bias in reject inference}
\author{Mateo Ordóñez Serna}
\date{20.08.2026}
\begin{document}
\maketitle
\begin{abstract}
The Reject Inference (RI) problem concerns the challenge of learning
from outcomes that are only observed for accepted applicants, and of
using information from rejected applicants whose outcomes remain
unknown. This selective observation induces acceptance bias and
motivates methods for its identification and correction. This thesis
introduces a dynamic perspective on RI by proposing surprise as a signal
of acceptance bias. \emph{Surprise} is defined as the discrepancy
between the expected performance of a classifier, as estimated from
historical data, and its realized performance over subsequent acceptance
rounds. The central hypothesis is that this discrepancy contains
information about the underlying acceptance bias and therefore provides
a means of observing it indirectly. First, the RI problem is formulated
within a measure-theoretic framework, making the underlying assumptions
explicit. Second, the hypothesis that surprise provides an observable
signal of acceptance bias is formally formulated. Third, a controlled
simulation study investigates the extent to which surprise reflects the
underlying bias. Finally, possible applications of the proposed signal
and directions for future research are discussed.
\end{abstract}

\renewcommand*\contentsname{Table of contents}
{
\setcounter{tocdepth}{3}
\tableofcontents
}

\newpage{}

\section{Introduction}\label{introduction}

The Reject Inference (RI) problem concerns learning from outcomes
observed only for accepted applicants while leveraging information from
rejected applicants whose outcomes remain unknown. This selective
observation induces acceptance bias and motivates methods for its
identification and correction. Reject Inference has been widely studied
in applications such as credit scoring, underwriting, and medical
decision systems, as well as from a more abstract perspective as a
missing data problem. Although both empirical and theoretical findings
have been mixed, its potential impact in high-stakes decision systems
continues to motivate the search for the ``holy grail'' to tackle down
the problem.

Rather than proposing a new correction method, this work revisits the RI
problem itself from a different perspective. Acceptance decisions are
repeatedly based on historically selected data, thereby perpetuating
acceptance bias. These accepted applicants subsequently become part of
the historical data used to make future decisions. RI is therefore
inherently a dynamic problem. Although several authors, such as
\textcite{kozdoi_et_al_BE_BASL}, generate simulation data using such an
acceptance loop, the dynamic structure of the loop itself has received
little attention as a potential source of information. The intuition
fueling this work is straightforward: acceptance bias likely operates as
a self-reinforcing mechanism. As a biased selection rule is repeatedly
applied, the resulting bias may compound over time or stabilize around a
sub-optimal equilibrium.

To study this intuition, a rigorous and transparent methodology is
indispensable. First, a short literature review is presented,
highlighting critiques and methodological shortcomings in the RI
literature that the dynamic perspective may help address. Second,
responding to the critique of \textcite{ehrhardt2021rejectReview}
regarding the lack of formalism in many RI studies, the credit scoring
setting is used to formulate a precise definition of the RI problem,
bridging the real world and Math. To do so, Measure theory provides a
framework to distinguish clearly between stochastic components and the
institution's deterministic decision rules. Third, a simulation
framework based on the well-known two-class Linear Discriminant Analysis
(LDA) setting is introduced to generate the data used throughout this
research. Fourth, \emph{surprise} is introduced as a heuristic quantity
intended to indirectly capture acceptance bias. Fifth, the bias of the
logistic regression model used in the acceptance loop and the
cross-validation (CV) based cut-off selection mechanism are related to
the surprise measurements to evaluate whether surprise provides an
approximation of acceptance bias. Finally, an outlook on possible
applications of surprise in RI methods is presented as motivation for
future research.

\newpage{}

\section{Related Work}\label{related-work}

RI can be formulated as a selective observation problem and studied
within the framework of missing-data theory. As such it also must handle
statistical questions like identification, sample selection, as well as
estimation and evaluation under bias. The review of
\textcite{ehrhardt2021rejectReview} revealed many structural problems in
existing RI-methods. Also, just like \textcite{HandHanley1993CanRIWork},
it concludes that for RI to perform (distinctly) better than using only
the accepts, additional identifying assumptions or additional
information beyond the observed accepted outcomes and applicant
covariates. Furthermore, in a spirit close to
\textcite{ehrhardt2021rejectReview}, I believe that being explicit about
the assumptions both on the data generating process studied, and the
methodology help to rigorously assess the feasibility, potential,
limitations and possible application fields of an RI approach.

Because the central concern of this thesis is not merely whether a
particular RI procedure performs well, but under which conditions
information about rejected applicants can be obtained from the available
data, the literature is reviewed along four dimensions:

\begin{itemize}
\tightlist
\item
  What information from the rejects does the RI-method exploit, and
  under what assumptions does this information identify the target of
  interest?
\item
  Are the assumptions concerning the data-generating process,
  observation mechanism, and methodology explicitly stated and
  theoretically/heuristically justified?
\item
  How do authors evaluate the validity, effectiveness, robustness, and
  practical applicability of their proposed approach, and what
  limitations do they identify?
\item
  To what extent does the methodology account for the temporal and
  feedback structure induced by repeated acceptance decisions?
\end{itemize}

All of this information is summarized in table
Table~\ref{tbl-lit-review} in the appendix. Thereby, when necessary the
assumptions or limitations (sometimes inherent of the methodology) were
explicitly stated rather than limiting to extract what the authors said
explicitly. First is clear that assumptions on the DGP and the
methodology are often ignored, which from my read also relates to a
weaker discussion of the assessment of the goodness and practicality of
the model. Second, papers using more information than only
cross-sectional features and the labels of the accepts tend to also
conclude having better performance. Finally, no paper has been found
that has formulated the RI problem from a dynamic perspective nor has it
been researched whether this perspective can be used to obtain important
learning signals.

\newpage{}

\section{Theoretical background and problem
setting}\label{theoretical-background-and-problem-setting}

Following the spirit of \textcite{ehrhardt2021rejectReview}, I begin by
formally presenting the Reject Inference (RI) Problem in Credit Risk as
understood and approached in this thesis, so that the hypothesis
formulation and the theoretical properties of the proposed solution can
be precisely formulated. First and also following the approach of
\textcite{ehrhardt2021rejectReview} a detailed description in natural
language of the biasing mechanism in FIs is done, so that each concept
used can be defined properly avoiding confusions. Second a measure
theoretic characterization of the biasing process through its assumed
underlying random variables (RV) is given. Third a precise formulation
of the missing-data mechanism guided by the definitions of
\textcite{littleAndRubingAnalysisWithMissingData} is presented, so that
the bias which RI attempts to alleviate can be clearly recognized
including its source. Finally, a notation of the available data for the
FI is done as observed realizations of the described processes. Although
the main focus of this thesis is not bettering the formalization
approaches, I find it important to be thorough so that future work can
properly criticize this thesis.

\subsection{Concepts definition}\label{sec-concepts-definition}

When a financial institution (FI) grants a credit, it has first received
information related to an applicant. This information consists of
features, some of which are (almost always) used in a scoring model
(``model features''), while others---e.g.~whether the applicant dresses
well in a face-to-face application---are not (``hidden features'').
Based on the available information, the FI decides whether to approve
the loan. Upon approval by the FI, a financing contract can be
concluded, leading to disbursement. A credit agreement after drawdown is
termed ``accepted''.\footnote{Since the contract requires the consent
  (``acceptance'') of both parties.} If an application does not lead to
a payout by the FI, it is called ``rejected''.\footnote{No distinction
  is made between rejection by the FI and withdrawal by the applicant.}

If an accepted loan is not repaid on time or in full (as defined by the
FI), it is classified as a default (``bad'' loan); otherwise, it is
classified as ``good''. Whether an applicant turns out to be a payer or
a non-payer is driven by their individual circumstances\footnote{e.g.~income
  or health} as well as purely unpredictable events\footnote{such as a
  bankruptcy due to a corruption scandal in the applicant's firm, or the
  applicant winning the lottery}, which I term idiosyncratic shocks. The
individual circumstances can be understood as a set of characteristics
consisting of the model features, hidden features and additional latent
factors, which cannot be observed by the FI. Ignoring idiosyncratic
shocks for a moment, the repayment outcome can be thought of as being
determined \emph{completely} by this full set of features. In this
sense, being a payer or non-payer can be viewed as an underlying
property of the applicant, even though it is not directly observable at
the time of the decision. Therefore, an applicant is called a \emph{good
applicant} if in the absence of an idiosyncratic shock, it would have
resulted in a good loan (if the loan was accepted). A \emph{bad
applicant} is defined analogously.\footnote{ Although the usage of bad
  and good for the loan and for the applicants in a slight different
  manner might feel confusing, it makes the mathematical notation
  \emph{much} easier.}

This interpretation does not imply that the FI could construct a perfect
prediction model. Even if all features---including hidden ones---were
observable, the functional relationship governing repayment would remain
unknown. In practice, both limited observability and unknown dependence
prevent perfect prediction.

The bias that reject inference addresses is not yet evident. To see it,
consider the acceptance process. At the beginning of its operations, the
FI has no data about its own clients. It must therefore rely either on a
scorecard trained on data from other markets or on rules based on expert
judgment. After sufficient data has been collected, the FI can train a
scorecard using its own historical data of \emph{accepted applicants},
since repayment outcomes are only observed for financed loans. As a
result, the model is estimated on a subset of the applicant
population---presumably one with higher creditworthiness---and then
applied to the full population of applicants. The distributional
differences between these populations, which have been well documented
(see \textcite{kozdoi_et_al_BE_BASL}), are (assumed) to be a central
source of bias highlighted in the reject inference literature.

Finally, note that the acceptance process is sequential. Newly accepted
applicants generate additional observations, which are incorporated into
the historical data and used to re-estimate the scorecard. This, in
turn, changes future acceptance decisions. In other words, the decision
rule evolves over time, and this time dimension must be considered when
modeling the reject inference problem. To capture this, I group time
into \emph{acceptance rounds}, where each round begins with a
re-estimation of the scorecard.

\subsection{Measure theoretic
description}\label{sec-measure-theoretic-desc-ri-problem}

Now with clear concepts from an economic perspective, we can define a
stochastic model for the RI problem. I try to be rigorous with the
formulation, since I believe it gives a language to precisely identify
the hidden assumptions most authors do when investigating RI. However,
in some cases I am rather loose on notation and formality -specially in
the notion as a stochastic process-, since I believe some nuances are
more confusing than helpful, particularly for a master's thesis.

Let \(\mathbf{T} = \{0, 1, ..., T\} \subseteq {\mathbb{N}}_0\) be the
set indexing the acceptance rounds and \(t \in \mathbf{T}\). Further let
\(X_m, X_h, X_l, Y, Y_m^{(t)}, Z^{(t)}, \epsilon\) and \(\varepsilon\)
be RVs on the probability space \((\Omega, \mathcal{F}, \mathbb{P})\).
Define with \(j \in \{m, h, l\}\) \[
X_j: (\Omega, \mathcal{F}) \to \left(\mathbb{R}^{f_j}, \mathcal{B}(\mathbb{R}^{f_j}) \right), \quad
f_j \in \mathbb{N}_0
\] as the RVs generating the applicant's model features, hidden features
and latent factors\footnote{ Note that even though not all features are
  per se continuous and therefore naturally real numbers, a reasonable
  assumption is that they are representable as real numbers}. Further
see that the case \(f_j = 0\)\footnote{Note: \(\mathbb{R}^{0} = \{()\}\)
  is a singleton set containing the empty tuple and
  \(\mathcal{B}\left(\mathbb{R}^{0}\right) = \{\emptyset, \{()\}\}\)}
represents a model where \(X_j\) always returns an empty tuple, meaning
that there are no features of type \(j\). I assume that at least
\(f_m > 0\). The joint observable feature representation is given by
\begin{equation}\protect\phantomsection\label{eq-X-map-definition}{
X : (\Omega, \mathcal{F}) \to \left(\mathbb{R}^{f}, \mathcal{B}\left(\mathbb{R}^{f} \right) \right), f:=f_m + f_h
}\end{equation} with
\begin{equation}\protect\phantomsection\label{eq-X-joint-observable-feats-vec}{
\forall \omega \in \Omega: X(\omega) := \begin{pmatrix} X_m(\omega) \\ X_h(\omega) \end{pmatrix}
}\end{equation} Further with \(V\) a placeholder for \(Y, Z^{(t)}\) and
\(Y_m^{(t)}\) let \[
V : (\Omega, \mathcal{F}) \to \left( E_V, 2^{E_V} \right)
\] with \[
E_Y = \{g,b\}, \quad E_{Z^{(t)}} = \{a,r\}, \quad E_{Y_m^{(t)}} = \{g,b, \texttt{na}\}
\]

be the RVs describing the repayment type, acceptance decision and
observable repayment outcome of the loan applied for defined
\(\forall \omega \in \Omega\) as \[
Y(\omega) = \begin{cases}
    b & \text{if applicant is a defaulter} \\
    g & \text{if applicant is a payer}
\end{cases} ,
\] \[
Z^{(t)}(\omega) = \begin{cases}
    a & \text{if applicant gets accepted (financed)} \\
    r & \text{if applicant gets rejected (non-financed)}
\end{cases}
\] and
\begin{equation}\protect\phantomsection\label{eq-y-m-t-as-measurable-function}{
Y_m^{(t)}(\omega) = \begin{cases}
    Y(\omega) & Z^{(t)}(\omega) = a \\
    \texttt{na} & Z^{(t)}(\omega) = r
\end{cases}
}\end{equation}

\(Z^{(t)}\) is more precisely modeled as a measurable function of both
features sets. Let \[
s_t : \mathbb{R}^{f_m} \to \mathbb{R}, \; 
\mathcal{B}\left(\mathbb{R}^{f_m}\right)-\mathcal{B}\left(\mathbb{R}\right) \text{ measurable}
\] be a score function and \[
d : \mathbb{R} \times \mathbb{R}^{f_h} \to \{a, r\}, \;
    \left(\mathcal{B}\left(\mathbb{R}\right)
    \otimes
    \mathcal{B}\left(\mathbb{R}^{f_h}\right)\right) - 2^{\{a, r\}}
    \text{ measurable}
\]

an acceptance decision function. Then

\begin{equation}\protect\phantomsection\label{eq-z-t-as-measurable-function}{
\forall \omega \in \Omega: \:
Z^{(t)}(\omega) =
d_t\left(
    s_t\left(X_m(\omega)\right),
    X_h(\omega)
\right)
}\end{equation}

In this spirit, \(Y\) can also be modeled as
\begin{equation}\protect\phantomsection\label{eq-y-as-measurable-function}{
\forall \omega \in \Omega:
Y(\omega) = \begin{cases}
    \varepsilon(\omega) & \epsilon(\omega) = \text{idiosyncratic shock} \\
    \iota \left(X_m(\omega), X_h(\omega), X_l(\omega) \right) & \epsilon(\omega) = \text{no idiosyncratic shock}
\end{cases}
}\end{equation}

with \(\iota\) a
\(\bigotimes_{j \in \{m, h, l\}} \mathcal{B}\left(\mathbb{R^{f_j}}\right)\)
-\(2^{E_Y}\) measurable function and \[
\epsilon : (\Omega, \mathcal{F}) \to \left(E_\epsilon, 2^{E_\epsilon} \right), \quad
\varepsilon : (\Omega, \mathcal{F}) \to \left(E_Y, 2^{E_Y} \right)
\]

the RVs defining whether an idiosyncratic shock occurs and its effect
(e.g.~if there is a shock in the labor market and whether it is positive
or negative for the applicant). From this perspective is easy to
precisely formulate the difference between a good/bad applicant and a
loan. A loan is a bad loan if \(Y = b\), an applicant is a bad applicant
if \(Y = b | \{\epsilon = \text{no external effect}\}\)

See that from Equation~\ref{eq-y-as-measurable-function} in connection
with Equation~\ref{eq-z-t-as-measurable-function} the FI is modeled to
be a ``good selector'' of features, since all the variables which it
uses to make the acceptance decision are assumed to actually play a role
in determining \(Y\). Although it might sound like an unrealistic
assumption, modeling variables unrelated to \(Y\) that affect
\(Z^{(t)}\) would render reject inference less useful or in some sense
easier, as the missingness mechanism would render ``closer to
ignorable''. An exact formulation of this idea is given in
Section~\ref{sec-formalization-missing-data-mechanism} after
Equation~\ref{eq-mnar-only-model}.

Now let the joint distribution of the applicants be given by
\begin{equation}\protect\phantomsection\label{eq-joint-dist-feats-and-repayment}{
(X_m, X_h, X_l, Y) \sim \mathbb{P}_{X_m, X_h, X_l, Y}
}\end{equation} defined over the product space of the state spaces of
the random variables, with marginals as above.\footnote{ Since all
  feature spaces are Borel spaces, their finite product is also a Borel
  space. As \(\left(E_Y, 2^{E_Y}\right)\) is finite, the resulting
  product space satisfies the regularity assumptions ensuring the
  existence of regular conditional probability kernels \autocite[Ch. 4,
  Theorem 9]{Pollard_2001}, which will later be required for conditional
  modeling arguments.} Furthermore I model the applicants' population at
each acceptance round \(t \in \mathbf{T}\) as an i.i.d. family \[
\left\{ \left(X_{m,i}^{(t)}, X_{h,i}^{(t)}, X_{l,i}^{(t)}, Y_{i}^{(t)} \right)\right\}_{i=1}^{N_t}
\] where each tuple of random variables in the family is distributed
according to \(\mathbb{P}_{X_m, X_h, X_l, Y}\). Meaning that the
features and repayment outcomes of applicants follow a time-invariant
distribution.

Therefore, the only source of temporal dynamics relies on the acceptance
decision. More specifically it arises from the estimation of \(s_t\). To
do so we also need to define the joint distribution of the observable
variables to the FI at each acceptance round as \[
\left(X_m, X_h, Y_m^{(t)} , Z^{(t)}\right) \sim \mathbb{P}_{X_m, X_h, Y_m^{(t)} , Z^{(t)}},
\] which also has marginals defined as above.\footnote{ Of course, the
  marginals of the \(X_m\) and \(X_h\) components correspond to the
  marginals according to \(\mathbb{P}_{X_m, X_h, X_l, Y}\).
  Argumentation for regularity is analogous to the former.} From this
defining the observable part of the applicant's distribution is
straightforward: \[
\left\{\left(X_{m,i}^{(t)}, X_{h,i}^{(t)}, Y_{m, i}^{(t)}, Z^{(t)}_i \right)\right\}_{i=1}^{N_t},
\] which, conditional on \(s_t\), forms an i.i.d. family for each fixed
acceptance round. This conditional i.i.d. characteristic already hints
that across acceptance rounds the distribution of each family may (and
likely does) change. These changes are not exogenous, but dependent on
the families of earlier acceptance rounds through estimation of \(s_t\).
Further consider that the probability kernels governing the feature
variables remain time invariant; the distributional shifts affect only
\(Z^{(t)}\) and therefore \(Y_{m}^{(t)}\).

This estimation of \(s_t\) can and should be further formalized to be
transparent on the further assumptions. Let \[
O_{:t} := \mathop{\Large\times}\limits_{\tau = 0}^{t-1} \mathop{\Large\times}\limits_{i=1}^{N_\tau}
        \left(\mathbb{R}^{f_m} \times E_{Y_m^{(\tau)}} \right)
\] with \(t \in \mathbf{T} \setminus \{0\}\) denote the product of the
underlying sets of all past model features and the observable repayment
outcomes. \(\mathcal{O}_{:t}\) denotes then the corresponding product
sigma algebra. For \(t=0\), \(s_t\) is assumed to come from some
external source (experts opinion, scoring trained on other data). For
\(t>0\), \(s_t\) is assumed to be generated through a measurable
learning procedure \[
\ell_t : \left(O_{:t}, \mathcal{O}_{:t}\right) \to \left(S_t, \mathcal{S}_t \right)
\] where the codomain denotes the measurable space of all admissible
score functions for time \(t\). See that \(\ell_t\) does not restrict
\(s_t\) to always be estimated by some specific procedure, such that the
score function at one acceptance round could be based on a XGB
classification forest and in another by a support vector machine.
Finally, it should be noted that although for \(d\) a similar definition
of some decision modeling procedure could be done, it is assumed in this
thesis that given \(s_t\), \(d\) does not change across \(t\), since
such dynamics are not considered in this thesis.

\subsection{Formalization of the missing data
mechanism}\label{sec-formalization-missing-data-mechanism}

Here the data missingness mechanism which gets recorded by \(Y^{(t)}_m\)
gets formalized according to the canonical definitions of
\textcite{littleAndRubingAnalysisWithMissingData}. Let in this section
\(t \in \mathbf{T} \setminus \{0\}\) be arbitrary but fixed, \[
X_m^{:t} = \bigcup_{\tau = 0}^{t-1}
 \left\{X^{(\tau)}_{m,i}\right\}_{i=1}^{N_\tau}
\] and \(X_h^{:t}\), \(X^{:t}\), \(Y^{:t}\), \(Y_m^{:t}\) be defined
analogously represent the ``flattened'' past random variables in the
processes.

The three canonical cases of missingness are Missing Completely at
Random (MCAR), Missing At Random (MAR) and Missing Not At Random (MNAR).
The first case means that the missing data is independent of covariates
and \emph{and} labels, meaning \[
\mathbb{P}\left(Z^{(t)} = r \middle | Y^{:t}, X_m^{:t}, X_h^{:t} \right) =
\mathbb{P}\left(Z^{(t)} = r \right)
\] so \(Z^{(t)}\) is independent of the three sets. This would be
observed if the FI rejected fully at random (e.g.~depending on the mood
of the employees), a ``through the door'' acceptance policy (always
accept) or would always reject. In any of these three cases, the scoring
function \(s_t\) would never be biased, meaning its parameters
\(\theta\) are expected to be calculated ``correctly''.

Here is important to explain what exactly biasing in this context means,
as talking about an ``unbiased estimator'' for methods like XGB is not
the same as for Maximum Likelihood estimators. Biasing is defined here
relatively loose: Given knowledge of the statistical properties of the
estimation method for \(s_t\) and enough data, one has good reasons to
believe that the estimated \(s_t\) will be a good approximation of some
scoring proportional to \(\mathbb{P}(Y = b | X)\). In this sense an
unbiased data-set occurs when it is representative enough (again,
assuming enough data, the small sample problem is not addressed here) to
get such a \(s_t\).

MAR on the other hand happens when the missingness only depends on the
covariates but not on the outcome variable. Here I define MAR as the
case where the missingness depends only on the model features, i.e.
\begin{equation}\protect\phantomsection\label{eq-mar-only-model}{
\mathbb{P}\left(Z^{(t)} = r \middle | Y^{:t}, X_m^{:t}, X_h^{:t} \right) =
\mathbb{P}\left(Z^{(t)} = r | X_m^{:t} \right)
}\end{equation} Note that even though this definition is conceptionally
the same as in many other sources, it should not be confused with
\begin{equation}\protect\phantomsection\label{eq-mar-all-x}{
\mathbb{P}\Big(
    Z^{(t)} = r \Big | Y^{:t},
    \underbrace{X_m^{:t}, X_h^{:t}}_{X^{:t}}
\Big)
=
\mathbb{P}\left(Z^{(t)} = r \middle | X_m^{:t}, X_h^{:t}\right),
}\end{equation} which could also be a valid definition of MAR but with a
different nuance.

Equation~\ref{eq-mar-all-x} is a valid definition of MAR in the sense
that a sample composed of \(Z^{:t}, Y^{:t}, X_m^{:t}, X_h^{:t}\) would
be MAR, and specially since for \(f_h=0\) it is equivalent to
Equation~\ref{eq-mar-only-model}. However, for \(f_h>0\) using a MAR
definition like Equation~\ref{eq-mar-all-x} would imply that a scoring
function \(s_t\) does not get trained on MAR data, as
\(Z^{(t)} | X_m^{:t} \, \cancel{\perp\!\!\!\perp} \, Y^{:t}\).
Considering that RI aims to reduce or eliminate the sample bias to get
an unbiased score function, it makes sense to focus on the bias that
does concern \(s_t\). Furhter, within the present framework,
Equation~\ref{eq-mar-all-x} follows directly form the definition of
\(Z^{(t)}\) -see Equation~\ref{eq-z-t-as-measurable-function}- as a
measurable function of \(X_m\) and \(X_h\).

If not only Equation~\ref{eq-mar-only-model} is fulfilled, but we also
have no omitted variable bias, i.e.~\(f_l=0\) (there are no latent
factors also determining \(Y\)) as well, Maximum Likelihood estimators
are also consistent under this kind of missingness mechanism
\autocite{littleAndRubingAnalysisWithMissingData}, even if the model
family is wrong, as Quasi-Maximum-Likelihood consistency results would
also apply, rendering ignorable\footnote{ Ignorable in the sense of
  \textcite{rubinMissingMechanismsOriginal}, i.e. one can train the
  model using only the data using only complete cases and still obtain
  consistent or unbiased results.} the missing data mechanism. Of course
if an estimator is consistent, then it also generates an unbiased
scoring function in the sense of this thesis. This argument is often
used by FIs and critics of RI to defend the usage of classical
statistical estimators like the logistic regression without needing to
consider the data missingness, either assuming \(f_l=0\) or that the
model features contain enough information of the latent factors, so that
the omitted variable bias is negligible. If the FI collects enough data
of the applicant, this might be a somewhat credible claim, but is it
really enough to just ignore the missingness mechanism? Later the
hypothesis checking section
(Section~\ref{sec-empirical-hypothesis-checking}) will give some
perspective on this regard.

The last case (MNAR) is the hardest case, as the missing data cannot be
ignored and still be able to train an unbiased scoring function.
Formally it can be written as a negation of
Equation~\ref{eq-mar-only-model}, i.e.
\begin{equation}\protect\phantomsection\label{eq-mnar-only-model}{
\mathbb{P}\left(Z^{(t)} = r \middle | Y^{:t}, X_m^{:t}, X_h^{:t} \right) \neq
\mathbb{P}\left(Z^{(t)} = r | X_m^{:t} \right)
}\end{equation} which can be intuitively understood as the acceptance
decision being (also) driven by factors not considered by \(s_t\) but
upon which \(Y^{:t}\) does depend on. This is characterized with
\(f_h > 0\). Note that the dependence of \(Y^{:t}\) to \(X_h^{:t}\) is
necessary for MNAR-ness. If one would use some factor independent of
\(Y^{:t}\) for the decision (e.g.~does the applicant have blond hair?),
decisions overwritten by this kind of factors would effectively render a
MCAR sample, since in these cases
\(Z^{(t)}_i  \, \perp\!\!\!\perp \, \left\{Y^{:t}, X_m^{:t}\right\}\).\footnote{
  Note: See that this definition of \(X_h\) as RVs on which \(Y\)
  depends upon is related to the case discussion of
  \textcite{bansikCrook2007} on the variables sets to be used for the
  acceptance-reject model and the default-repayment model. Such a
  discussion is however not necessary here given the probabilistic model
  assumed in this thesis.} Additionaly see that this discussion makes it
natural to alternatively (and equivalently) write the MNAR condition as
\begin{equation}\protect\phantomsection\label{eq-mnar-as-not-independence}{
Z^{(t)} | X_m^{:t} \cancel{\perp\!\!\!\perp} Y^{:t}
}\end{equation}

Furthermore consider that MNAR (at least as modeled here) inherently
implies an omitted variable bias, as
\(Y^{:t} \cancel{\perp\!\!\!\perp} X_h^{:t}\). This makes the evaluation
of the effectivness of RI as a tool to correct for MNAR complicated, as
one cannot cleanly separate the effect of MNAR-ness and of the omitted
variable on the performance loss of the score-card estimation. At most
one can reduce the effect of the omitted variable bias by increasing the
dependence structure between \(X_m\) and \(X_h\), however this would on
the other hand \emph{also} weaken the dependence described by
Equation~\ref{eq-mnar-as-not-independence} and in this sense approach
MNAR-ness to the MAR case.

\newpage{}

\section{Simulation framework}\label{simulation-framework}

This thesis uses an acceptance loop which has been further developed
from the loop described in \textcite{kozdoi_et_al_BE_BASL}. In this
section the components of the loop are defined, i.e.

\begin{itemize}
\tightlist
\item
  The joint distribution \(\mathbb{P}_{(Y, X)}\) of the repayment
  variable and the covariates,
\item
  the acceptance decision \(Z^{(t)}\) for each acceptance round, meaning
  formulating

  \begin{itemize}
  \tightlist
  \item
    the estimation method of the score function \(s_t\),
  \item
    the threshold calculation and
  \item
    the complete decision function \(d\) for both the MAR and MNAR cases
  \end{itemize}
\item
  from which the rest of the RVs and functions get formulated.
\end{itemize}

After defining each of the components an overview of the process as an
algorithm is provided.

\subsection{Applicants' population}\label{sec-app-population-simulation}

Recalling section Section~\ref{sec-measure-theoretic-desc-ri-problem},
according to the probabilistic model for the RL problem it suffices to
define the distributions of the repayment outcome and the features to
define the applicants' population. In the simulation I do not consider
latent features (case \(f_l = 0\)). Further let
\(\tilde{Y} := Y | \{\epsilon = \text{no external effect}\}\) denote the
``noiseless'' repayment outcome. I describe the population with this
variable, since it is less convoluted, is the object that the scorecard
will attempt to model and is easily bridged to the idea of a good and
bad applicants. The clean model corresponds to the classical two class
multivariate Gaussian problem of linear discriminant analysis. It might
feel out of the decade to use such a simple DGP, however using the very
well known statistical properties of this model allows carefully study
the RI framework presented here very carefully.

It holds \(\tilde{Y} \sim Bernoulli(\pi_{b})\) with \(\pi_b \in (0,1)\)
the probability of a bad applicant. The observable features vector
\(X\)\footnote{ Recall from
  Equation~\ref{eq-X-joint-observable-feats-vec}:
  \(X = \left(X_m^{\top}, X_h^{\top} \right)^{\top}\)} is defined as a
Gaussian Mixture with two components, a good one (paying population) and
a bad one (defaulters). With \(o \in \{g, b\}\) we can characterize the
distributions of these populations as \[
X | \{\tilde{Y}=o\} \sim \mathcal{N}\left(
    \mu_o =
    \begin{pmatrix}
        \mu_{m, o} \\ \mu_{h, o}
    \end{pmatrix},
    \Sigma_o =
    \begin{pmatrix}
        \Sigma_{m, o} & \Sigma_{mh, o} \\
        \Sigma_{mh, o}^{\top} & \Sigma_{h, o}
    \end{pmatrix}
\right)
\] with \(\Sigma_o, \Sigma_{m, o}\) and \(\Sigma_{h, o}\) strictly
positive definite (so, non-singular). Letting \(X_o\) be short for
\(X | \{\tilde{Y}=o\}\) also let \(\phi_{X_o}\) denote its probability
density function (pdf). Since \(X_o\) is normally distributed the pdf is
given by \[
\phi_{X_o}(x) := 
\frac{1}{\sqrt{(2\pi)^f |\Sigma_o|} }
\exp\left\{-\frac{1}{2} (x-\mu_o)^\top \Sigma_o^{-1} (x-\mu_o)\right\}
\]

For further analysis it is also helpful to write down the
characterization of the joint distribution of \((X,\tilde{Y})\), which
is the marginal distribution of \(\mathbb{P}_{X_m, X_h, X_l, Y}\) over
the latent feature component, given
\(\epsilon = \text{no external effect}\). Its density\footnote{
  Following \textcite[Ch. 3, Sec. 1]{Pollard_2001}, densities are
  understood here in the Radon--Nikodym sense, such that absolute
  continuity with respect to Lebesgue measure alone is not required, but
  can be given with respect to any dominating measure.} is defined as a
Radon-Nikodym derivative, using that the normal distribution is
absolutely continuous w. r. t. the Lebesgue Measure, in this case
\(\lambda^f\) over \(\mathbb{R}^f\), and the discreteness of \(Y\) \[
p_{X, \tilde{Y}} = \frac{d\mathbb{P}_{X, \tilde{Y}}}{d\left(\lambda^f \otimes \#_Y\right)}
\] where \(\#_Y\) is the counting measure over \(E_Y\), which gets
evaluated pointwise as a map \(\mathbb{R}^f \times E_y \to [0, \infty)\)
\begin{equation}\protect\phantomsection\label{eq-join-density-x-tilde-y}{
p_{X, \tilde{Y}}(x, o) = \pi_o \cdot \phi_{X_o}(x)
}\end{equation} where \(\pi_g := 1 - \pi_b\).

From this we can easily derive the marginals\footnote{ Since
  \(\pi_b \in (0, 1)\) and \(\phi_{X_o}\) is strictly positive
  \emph{and} finite everywhere, the marginal densities of
  \(\mathbb{P}_{X, \tilde{Y}}\) over \(X\) and \(\tilde{Y}\) each are
  well-defined \autocite[Ch. 4, Theorem 12]{Pollard_2001}} as
\begin{equation}\protect\phantomsection\label{eq-marginal-of-Y}{
p_{\tilde{Y}}(o) = \int_{\mathbb{R}^{f}} p_{X, \tilde{Y}}(x, o) \, \lambda^f(dx) = \pi_o ,
}\end{equation} which of course corresponds to the probability mass
function of a Bernoulli distributed variable, and
\begin{equation}\protect\phantomsection\label{eq-marginal-of-X}{
p_X(x) = \int_{E_Y} p_{X, \tilde{Y}}(x, y) \, \#_{Y}(dy) = \pi_b \phi_{X_b}(x) + \pi_g \phi_{X_g}(x)
}\end{equation}

which is the pdf of a Gaussian Mixture and thus finally:
\begin{equation}\protect\phantomsection\label{eq-posterior-of-Y}{
\mathbb{P}(\tilde{Y} = o | X=x) = p_{\tilde{Y} | X} (o, x) = \frac{\phi_{X_o}(x) p_{\tilde{Y}}(o)}{p_X(x)}
}\end{equation}

Now let us assume that \(\epsilon, \varepsilon\) are Bernoulli
distributed with parameters
\(\pi_{\epsilon} := \mathbb{P}(\epsilon =\text{no idiosyncratic shock}) > 0\)
and \(\pi_{\varepsilon, b} := \mathbb{P}(\varepsilon = b) \in (0,1)\).
With this we can further write down
\begin{equation}\protect\phantomsection\label{eq-pd-with-shock-perf-bayes}{
\mathbb{P}(Y = b | X = x) = p_{Y=b | X}(x) = 
    \pi_{\epsilon} \, p_{\tilde{Y} | X} (b, x) +
    (1 - \pi_{\epsilon}) \pi_{\varepsilon, b}
}\end{equation} which is the basis for the best possible prediction of a
bad loan.

However, since the FI ignores \(X_m\) for the scorecard creation,
\(s_t\) in reality can only approximate
\(\mathbb{P}(Y = b | X_m = x_m)\). This distribution can also get
quickly calculated similarly by integrating out \(X_h\) from the steps
above, see for example that we can write
\(x = \left(x_m^\top, x_h^\top \right)^\top\) in
Equation~\ref{eq-join-density-x-tilde-y}, such that \[
p_{X_m, \tilde{Y}} (x_m, o) :=  \int_{\mathbb{R}^{f_h}} p_{X, \tilde{Y}}(x, o) \lambda^{f_h}(d x_h)
\]

from that it follows that \[
X_{m,o} := X_m | \tilde{Y} = o \sim \mathcal{N} \left(\mu_{m, o}, \Sigma_{m, o} \right)
\] which is just the RV induced by a marginal distribution of \(X\).
Then we can calculate that
\begin{equation}\protect\phantomsection\label{eq-pd-with-shock-perf-bayes-reduced-model}{
\mathbb{P}(Y = b | X_m = x_m) = p_{Y=b | X_m}(x_m) = 
    \pi_{\epsilon} \,
    \underbrace{\frac{\phi_{X_{m, b}}(x_m) p_{\tilde{Y}}(b)}{p_{X_m}(x_m)}}_{
        =: p_{\tilde{Y} | X_m} (b, x_m)
    }
    +
    (1 - \pi_{\epsilon}) \pi_{\varepsilon, b}
}\end{equation}

which does represent the best estimation we can hope to get from
\(s_t\).

For the simulation I have chosen a fairly simple DGP, so that its
statistical properties are easily understood. Namely, \(f_m = 2\),
\(f_h = 1\), \(\Sigma_g = \Sigma_b = \Sigma\), such that the
distributional differences between \(X_o\) and \(X_g\) are solely
characterized by the mean difference. The covariance matrix is
characterized as follows
\begin{equation}\protect\phantomsection\label{eq-general-form-Sigma}{
\Sigma = \begin{pmatrix}
1 & \sigma_{m, 12} & \sigma_{mh} \\
\sigma_{m, 12} & 1 & \sigma_{mh} \\
\sigma_{mh} & \sigma_{mh} & 1
\end{pmatrix}
}\end{equation}

and the means are given by

\[
\mu_{b} = \begin{bmatrix}
    0.0\\
    0.0\\
    0.0
\end{bmatrix}, \; \mu_{g} = \begin{bmatrix}
    1.0\\
    2.0\\
    -1.0
\end{bmatrix}.
\]

The form in Equation~\ref{eq-general-form-Sigma} for the variance
covariance matrix was selected since it allows to study the full space
of positive definite variance covariance matrices in this form in order
to choose sensible values for \(\sigma_{m, 12}\) and \(\sigma_{mh}\). A
useful way is to look at the performance of the classifiers given by
Equation~\ref{eq-pd-with-shock-perf-bayes} and
Equation~\ref{eq-pd-with-shock-perf-bayes-reduced-model}, which are the
bayes classifiers\footnote{ Also known as ``the generative model'' in
  the linear discriminant analysis literature} for \(Y | X\) and
\(Y | X_m\) respectively. Their (theoretical) performance can be
measured as their (Bayes) error rate, which is the probability that the
classifier makes a wrong classification, given that the decision rule is
to predict bad if the predicted probability of bad is higher than the
predicted probability of good and otherwise predicts good.\footnote{ See
  that this rules does not allow for the option ``do not predict
  anything''} This probability can be analyitically calculated with
Equation~\ref{eq-bayes-error-rate-formula} as derived in the appendix
(see Section~\ref{sec-derivation-bayes-rate}). Let \(P_e^{*}\) denote
the error rate of the full model - so following
Equation~\ref{eq-pd-with-shock-perf-bayes}- and \(P^m_e\) the error rate
of the reduced model according to
Equation~\ref{eq-pd-with-shock-perf-bayes-reduced-model}.

In general it holds \(0 \leq P_e^{*} \leq P^m_e\) and because of the
properties of Gaussian distributions, \(P^m_e\) does not depend on
\(\sigma_h\). Using this properties we can check the distribution of
these error rates for all admissible \(\Sigma\) matrices, namely all
positive definite ones. This set can also be analitically calculated
using Equation~\ref{eq-positive-definite-region} (see
Section~\ref{sec-set-pos-definite-vcovs} in the appendix for the
derivation). In Figure~\ref{fig-all-bayes-errs} the results of this
analysis is shown.

\begin{figure}[H]

\centering{

\pandocbounded{\includegraphics[keepaspectratio]{thesis_files/figure-pdf/fig-all-bayes-errs-output-1.pdf}}

}

\caption{\label{fig-all-bayes-errs}Representation of joint distribution
\((Y, X)\)}

\end{figure}%

In the left panel we see the heat map of the ratio \(P_e^{*}/P^m_e\).
The graphic shows numerically that
\begin{equation}\protect\phantomsection\label{eq-equal-error-rates}{
\forall \sigma_{m, 12} \exists \sigma_{mh} : P_e^{*}=P^m_e
}\end{equation}

and for many of the cases it does not require the error rate to be zero.
Since as it will be further explained in
Section~\ref{sec-hidden-feats-overwrites}, the MNAR-causing acceptance
decision mechanism depends entirely on \(X_h\), the cases in
Equation~\ref{eq-equal-error-rates} result very interesting, as they
allow to study MNAR-ness with almost no effect of the hidden variable
bias.

The last claim is non-trivial and should be argued for. A DGP with
\(P_e^{*}=P^m_e\) implies that for this DGP there exists a classifier
depending solely on \(X_m\) that has a discriminatory power as good as a
classifier using the full features vector. Further, see that the joint
distribution \((Y, X_m)\)

This set is specially interesting for studying the MNAR mechanism, which
as defined in Equation~\ref{eq-mnar-only-model} where the MNAR mechanism
is dependent on the latent factor \(X_h\).

Its Bayes Error rate is 14.78\(\%\). We can see a visualization of the
complete dgp in Figure~\ref{fig-3d-base-dgp}, whereby the Perfect Bayes
Boundary corresponds to the classification boundary of a classifier
given by Equation~\ref{eq-pd-with-shock-perf-bayes}. This boundary is
calculated in the appendix (see
Section~\ref{sec-dec-boundary-perf-bayes-idyiosincratic-shocks}) and can
be rewritten from Equation~\ref{eq-decision-boundary-equal-covs} as the
hyperplane
\begin{equation}\protect\phantomsection\label{eq-hyperplane-decision-boundary}{
W x - \beta = \log c
}\end{equation} with
\(W=(\mu_b - \mu_g)^\top \Sigma^{-1} = \Sigma^{-1}(\mu_b - \mu_g)\),
\(\beta = W\frac{\mu_b + \mu_g}{2}\) and \(c\) a real constant depending
on the idyiosincratic schock dynamics (\(\pi_{\epsilon}\) and
\(\pi_{\varepsilon, b}\)) and the unconditional probability of default
(\(\pi_b\)).

\begin{figure}[H]

\centering{

\pandocbounded{\includegraphics[keepaspectratio]{thesis_files/figure-pdf/fig-3d-base-dgp-output-1.pdf}}

}

\caption{\label{fig-3d-base-dgp}Representation of joint distribution
\((Y, X)\)}

\end{figure}%

Considering the bayes error rate of the problem, the good separability
by the hyperplane of the perfect bayes decision boundary observable in
Figure~\ref{fig-3d-base-dgp}, it is clear that a logit model (which also
has a linear discriminant) with no missing variable bias, should have no
problem estimating a good classifier. To check out how difficult the
problem gets while ignoring one of the variables, a pairwise component
plot as in Figure~\ref{fig-pairwise-base-dgp} is a good diagnostic.
There the perfect bayes boundary was calculated using the hyperplane
described in Equation~\ref{eq-hyperplane-decision-boundary}, whereby the
non considered component of the features in each plot was set to its
expected value of \(X\) given no shock. Meaning with \(r\) component's
position in \(X\) of the non-shown component, it was set in
Equation~\ref{eq-hyperplane-decision-boundary} \[
x\mathtt{[r]} = \mathbb{E}(X | \epsilon = \text{no shock})\mathtt{[r]} = (\pi_b)\mu_b\mathtt{[r]} + (1-\pi_b)\mu_g\mathtt{[r]}
\]

It is clear that Figure~\ref{fig-pairwise-base-dgp} also suggests that
the classification problem shoud also be well solved by a logistic
regression considering only two of the components of \(X\). This figure
also shows a bit the importance of \(X_h\) for the classification
problem in relationship to the correlation to \(X_{m}\). See that the
classification boundary in the graphs containing \(X_h\) has a positive
slope w. r. t. \(X_{m,1}\) and \(X_{m, 2}\). This means that a higher
(positive) correlation of the hidden factor to the model variables
increases the separation power when observing \(X_h\). A negative
correlation should approximately have the opposite effect.

\begin{figure}[H]

\centering{

\pandocbounded{\includegraphics[keepaspectratio]{thesis_files/figure-pdf/fig-pairwise-base-dgp-output-1.pdf}}

}

\caption{\label{fig-pairwise-base-dgp}Pairwise covariates component
representation of \(Y, X\)}

\end{figure}%

\subsection{The acceptance loop}\label{the-acceptance-loop}

To put in perspective the acceptance loop

Algorithm \ref{alg:ac-forward-pass}

\begin{algorithm} 
\caption{Actor-Critic Forward Pass}
\label{alg:ac-forward-pass}
\begin{algorithmic}[1] 
    
    \Procedure{ActorCritic}{$A_{t-1}, M_{t-1}$}
        \State $\mathbf{P} \gets [\ ]$, $A_t \gets [\ ]$ \Comment{$\mathbf{P}$roabilities and new $A$rchitecture lists}
        \State $H_t, M_t \gets \texttt{CriticLSTM}\left(A_{t-1}, \texttt{hx=}M_{t-1}\right)$
        \State $h_t \gets H_t\texttt{[-1, :]}$
        \State $\mathbb{P}_{N_{TSF}} \gets \operatorname{softmax}\left(\texttt{AH}_1^{MLP}\left(h_t\right)\right)$ \Comment{$\mathbb{P}_{N_{TSF}} \in \mathbb{R}^{\max \mathcal{N}_{TSF}}$}
        \State $N_{TSF} \gets \operatorname{MultiNomialSample}\left(\mathbb{P}_{N_{TSF}}\right)$
        \State $A_t\texttt{.append}\left(N_{TSF}\right)$
        \State $\mathbf{P}\texttt{.append}\left(\mathbb{P}_{N_{TSF}}\right)$
        \State $\text{LogProbSum} \gets \log\mathbb{P}_{N_{TSF}}\texttt{[}N_{TSF} \texttt{]}$
        \For{$\texttt{i}_A \in \{1, \ldots , N_{TSF}\}$}
            \State $\mathbb{P}_{\texttt{i}_A} \gets \operatorname{softmax}\left(W_t^{\texttt{i}_A, h}h_t + W_t^{\texttt{i}_A, p}\mathbf{P}\texttt{[-1]} + b_t^{\texttt{i}_A}\right)$ \Comment{$\mathbb{P}_{\texttt{i}_A} \in \mathbb{R}^{\#TSF}$}
            \State $TSF_{\texttt{i}_A} \gets 
            \operatorname{MultiNomialSample}\left(\mathbb{P}_{\texttt{i}_A}\right)$
            \State $\mathbf{P}\texttt{.append}\left(\mathbb{P}_{\texttt{i}_A}\right)$
            \State $A_t\texttt{.append}\left(TSF_{\texttt{i}_A}\right)$
            \State $\text{LogProbSum} \gets \text{LogProbSum} +  \log\mathbb{P}_{\texttt{i}_A}\texttt{[}TSF^e_{\texttt{i}_A} \texttt{]}$
        \EndFor
        \For{$\texttt{i}_A \in  \{1, \ldots , \max\mathcal{N}_{TSF}\} \setminus  \{1, \ldots , N_{TSF}\}$}
            \State $\operatorname{NoGrad}\left(W_t^{\texttt{i}_A, h}, W_t^{\texttt{i}_A, p}, b_t^{\texttt{i}_A}\right)$
        \EndFor
        \State $\mathbb{P}_{AG} \gets \operatorname{softmax}\left(W_t^{AG, h}h_t + W_t^{AG, p}\mathbb{P}_{N_{TSF}} + b_t^{AG}\right)$ \Comment{$\mathbb{P}_{AG} \in \mathbb{R}^{\#AG}$}
        \State $AG \gets \operatorname{MultiNomialSample}\left(\mathbb{P}_{AG}\right)$
        \State $\mathbf{P}\texttt{.append}\left(\mathbb{P}_{AG}\right)$
        \State $A_t\texttt{.append}\left(AG\right)$
        \State $\text{LogProbSum} \gets \text{LogProbSum} +  \log\mathbb{P}_{AG}\texttt{[}AG \texttt{]}$
        \State $V_{h_{t}} \gets \texttt{CriticPerceptron}\left(h_{t}\right)$ \Comment{Evaluate $h_{t}$ with Value Function}
        \State $\operatorname{SaveInClass}\left(\text{LogProbSum}, V_{h_{t}}\right)$
        \State \Return $\mathbf{P},\, A_t,\, M_t$ 
    \EndProcedure
\end{algorithmic} 
\end{algorithm}

\subsection{Decision and scoring
functions}\label{decision-and-scoring-functions}

At the beginning of the loop a business rule is used, where values of
the second component

The decision function has two phases, the starting phase (where no data
has been gathered yet) and the scoring function estimating phase. In the
following is explained how in each moment the function looks like.

\subsubsection{Starting the loop: business
rule}\label{starting-the-loop-business-rule}

Since it is hard to

\subsubsection{Scoring function
estimation}\label{scoring-function-estimation}

\subsubsection{CV based threshold}\label{sec-cv-based-threshold}

Then at each acceptance round, the chosen threshold by CV is denoted as
\(c_t\) and the decision function is therefore

\subsubsection{Hidden features
overwrites}\label{sec-hidden-feats-overwrites}

\begin{equation}\protect\phantomsection\label{eq-Z-t-with-threshold-and-scorecard-no-overwrites}{
Z^{(t)}_i(\omega) = \begin{cases}
a & \hat{s}_t \left(X^{(t)}_{m, i} (\omega) \right) > c_t  \\
r & \hat{s}_t \left(X^{(t)}_{m, i} (\omega) \right) \leq c_t
\end{cases}
}\end{equation}

\subsection{Acceptance loop algorithm}\label{acceptance-loop-algorithm}

\newpage{}

\section{Using the acceptance
dynamics}\label{using-the-acceptance-dynamics}

With a defined simulation framework, one can begin to formulate
precisely how to use the acceptance dynamics. The core idea is that if
the bias exists or is economically relevant it should be observable in
some manner and if so we can try to minimize the bias. Based on this
hypothesis one can then formulate a scenario, using that we have full
knowledge of the applicants' population, which allows gathering evidence
on the plausibility of such hypothesis. Having delivered this evidence,
a reward signal is defined along to a formulation of the RI problem as a
reinforcement learning problem.

\subsection{Hypothesis definition}\label{hypothesis-definition}

To formalize the statement ``the bias should be observable and can be
minimized'' we need to first formulate what is understood as bias and
how one can observe it. First I define what I understand as observable
bias, which I call surprise. Second, a theoretical motivation is
presented of when this bias should \emph{not} be observable, that gives
an intuition of the plausibility of the observability of the bias.
Third, I give some simulation based evidence of my claims on
observability of the bias.

\subsubsection{Surprise signal}\label{surprise-signal}

It is well known that biased samples lead to biased evaluation of
classifiers in the sense that is the performance is not representative
on how the classifier performs on unseen data
\autocite{sugiyama2007covariate}, as sampling bias results in covariate
shifts. Most of the literature on this topic has a focus on the bias as
measured on an unseen unbiased sample. However, I believe that this bias
is also observed on the population of accepted applicants, and it can be
used as an informative signal.

My claim is that this bias is also observable on the accepted applicants
as well. This means that the expected performance of a scorecard as
evaluated on accepts does not correspond to the observed performance on
accepted applicants of the next round. Then, I define the surprise
\(\eta^{(t)}\) at acceptance round \(t\) analogously to the definition
of evaluation bias, namely as the difference between expected
performance \(m_{CV}\) measured by a metric \(m\) generated through CV
on historical data and the realized performance on the new unseen data
of the accepts. \[
\eta^{(t)} := m_{CV}\left(\mathbf{X}_m^{:t}, \mathbf{Y}^{:t}\right) - m\left(\mathbf{X}_m^{(t)}, \mathbf{Y}^{(t)}\right)
\]

\subsubsection{An unsurprising
setting}\label{sec-an-unsurprising-setting}

The claim that the bias can be observed only through the accepted
applicants is strong and should therefore be motivated. The purpose of
this section is not to provide a rigorous proof, but rather an intuition
of when one should \emph{not} expect to observe surprise.

Consider a FI operating under the following assumptions:

\begin{enumerate}
\def\labelenumi{\arabic{enumi}.}
\tightlist
\item
  The applicants population distribution is time invariant.
\item
  The number of applicants \(N_t\) is sufficiently large at each round,
  so that many applicants are accepted.
\item
  The scorecard estimation algorithm remains unchanged across rounds,
  i.e.~\(\ell_t = \ell_{t'}\) for \(t, t' \in \mathbf{T}\).
\item
  The cutoff \(c_t\) is chosen consistently across rounds according to
  the same business criterion.
\item
  The missing data process satisfies the MAR assumption, meaning
  \(f_h = 0\) and \(Z^{(t)}_i\) is fully determined by \(\hat{s}_t\),
  \(c_t\) and \(X^{(t)}_{m,i}\) for
  \(t \in \mathbf{T}, i \in \{1, \ldots, N_t\}\).\footnote{ Recall
    Equation~\ref{eq-Z-t-with-threshold-and-scorecard-no-overwrites}.}
\end{enumerate}

Under MAR, the repayment behaviour of an applicant is independent of the
acceptance decision once the observable characteristics are known.
Therefore,

\[
Z^{(t)} \perp\!\!\!\perp Y^{:t} \mid X^{:t}
\]

which implies that the conditional relationship between repayment and
features remains unchanged by the acceptance process. Together with a
stationary applicants population and a stable estimation procedure, one
would expect the distribution of accepted applicants to remain
approximately stable across rounds.

In such a setting, the scorecard is repeatedly trained and evaluated on
samples generated by essentially the same acceptance mechanism.
Consequently, the performance estimated through cross-validation on
historical accepts should be close to the realized performance on future
accepted applicants. In other words, one would expect only small values
of \(\eta^{(t)}\)

This observation motivates the interpretation of surprise as a
diagnostic signal. Persistent or systematically large values of
\(\eta^{(t)}\) suggest that at least one of the assumptions above is
violated. Although the assumptions may be restrictive, assumptions 1 to
4 are enforced by the simulation framework. Consequently, within the
simulated environment persistent surprise cannot be attributed to
changes in the applicants' population, the estimation procedure or the
acceptance policy. Instead, surprise is expected to primarily reflect
model misspecification and departures from MAR. Therefore, at least
within the context of this thesis, surprise provides a natural candidate
reward signal for learning policies that mitigate the effects of such
departures.

\subsubsection{Empirical evidence of the bias observability as
surprise}\label{empirical-evidence-of-the-bias-observability-as-surprise}

Lets see

\begin{figure}[H]

\centering{

\pandocbounded{\includegraphics[keepaspectratio]{thesis_files/figure-pdf/fig-expectation-vs-observed-performance-output-1.pdf}}

}

\caption{\label{fig-expectation-vs-observed-performance}Performance on
the expectation samples and realized samples}

\end{figure}%

\begin{figure}[H]

\centering{

\pandocbounded{\includegraphics[keepaspectratio]{thesis_files/figure-pdf/fig-abs-difference-between-exp-and-obseved-performance-output-1.pdf}}

}

\caption{\label{fig-abs-difference-between-exp-and-obseved-performance}Absolute
difference between the performance on the expectation samples and
realized samples \((W_{MA}=5)\)}

\end{figure}%

\begin{figure}[H]

\centering{

\pandocbounded{\includegraphics[keepaspectratio]{thesis_files/figure-pdf/fig-areas-diff-logit-in-grid-output-1.pdf}}

}

\caption{\label{fig-areas-diff-logit-in-grid}Absolute difference between
the performance on the expectation samples and realized samples
\((W_{MA}=5)\)}

\end{figure}%

\begin{figure}[H]

\centering{

\pandocbounded{\includegraphics[keepaspectratio]{thesis_files/figure-pdf/fig-difference-between-exp-and-obseved-performance-output-1.pdf}}

}

\caption{\label{fig-difference-between-exp-and-obseved-performance}Difference
between the performance on the expectation samples and realized samples
\((W_{MA}=5)\)}

\end{figure}%

NEED TO FIND A WAY TO MAKE THIS SMALLER

\begin{figure}[H]

\centering{

\pandocbounded{\includegraphics[keepaspectratio]{thesis_files/figure-pdf/fig-diminishing-accepts-output-1.pdf}}

}

\caption{\label{fig-diminishing-accepts}Oracle's vs.~Biased model Accept
rates per generation round of base simulation}

\end{figure}%

\subsubsection{Formulating the
hypothesis}\label{formulating-the-hypothesis}

\subsubsection{Checking the hypothesis}\label{checking-the-hypothesis}

Although a theoretical check of the hypothesis might result interesting,
it is neither straigh

\paragraph{Classifiers definition}\label{classifiers-definition}

\paragraph{Evidence of hypothesis plausibility given the simulation grid
results}\label{sec-empirical-hypothesis-checking}

\subsection{RI as a dynamic RL
problem}\label{ri-as-a-dynamic-rl-problem}

\subsubsection{Reward signal definition}\label{reward-signal-definition}

\subsubsection{Problem formulation}\label{problem-formulation}

\subsection{RL Agent}\label{rl-agent}

\subsubsection{Explaining the
Data-Structure}\label{explaining-the-data-structure}

\subsubsection{RL-method used}\label{rl-method-used}

\paragraph{Why this method?}\label{why-this-method}

\paragraph{Presentation of the
architecture}\label{presentation-of-the-architecture}

\paragraph{Training loop algorithm}\label{training-loop-algorithm}

\paragraph{Transparency on
metaparameters}\label{transparency-on-metaparameters}

\newpage{}

\section{Results of Agent}\label{results-of-agent}

\subsection{``Base Experiment''}\label{base-experiment}

\subsection{Sensitivity analysis}\label{sensitivity-analysis}

\newpage{}

\section{Conclussion}\label{conclussion}

\subsection{Hypotheses fulfilled?}\label{hypotheses-fulfilled}

\subsection{Limitations of the study}\label{limitations-of-the-study}

\subsection{Future work (building on the limitations and highlighting
measure theoretic
possibilities)}\label{future-work-building-on-the-limitations-and-highlighting-measure-theoretic-possibilities}

\newpage{}

\section{Appendix}\label{appendix}

\subsection{Acknowledgement of AI
usage}\label{acknowledgement-of-ai-usage}

Generative Artificial Intelligence were used as interactive assistants
for discussing statistical concepts, reviewing mathematical arguments,
suggesting improvements to exposition, identifying relevant literature,
reviewing software implementations and help in generating code. The used
AI's were

\begin{itemize}
\tightlist
\item
  ChatGPT \autocite{openai2025chatgpt}, mostly GPT 5.1
\item
  Claude \autocite{claude}, mostly Sonnet 4.6
\item
  Microsoft Copilot \autocite{microsoft_copilot}, Microsoft Prometheus
  model powered by ChatGPT
\item
  Perplexit \autocite{perplexity}
\item
  GitHub Copilot \autocite{github_copilot}
\item
  DeepAI \autocite{deepai}
\end{itemize}

A close bookkeeping of the usage can be seen on
\href{https://github.com/TheBISEconStatcian/BEReBASL/blob/main/ai_usage/ai_usage.md}{GitHub}

\subsection{Mathematical
derivations/proofs}\label{mathematical-derivationsproofs}

\subsubsection{Set of positive definite variance covariance matrices for
the simulation}\label{sec-set-pos-definite-vcovs}

As stated in Section~\ref{sec-app-population-simulation}, the common
covariance matrix \(\Sigma\) has the general shape (for the 3 variables
case)

\begin{equation}\protect\phantomsection\label{eq-general-Sigma-3d}{
\Sigma =
\begin{pmatrix}
1 & \sigma_{m,1,2} & \sigma_{mh} \\
\sigma_{m,1,2} & 1 & \sigma_{mh} \\
\sigma_{mh} & \sigma_{mh} & 1
\end{pmatrix}
}\end{equation}

For notational simplicity let here \(a \equiv \sigma_{m,1,2}\) and
\(b \equiv \sigma_{mh}\).\footnote{ Note: from here one the derivation
  has been fully AI-generated with minor changes by me. The facts were
  checked upon by me. See
  \textcite{chatgpt_positive_definite_covariance_2026} for the reference
  to the chat.} The space of positive definite covariance matrices
following Equation~\ref{eq-general-Sigma-3d} can be found using
Sylvester's criterion\footnote{ I.e. if all leading principal minors
  have a positive determinant, then the matrix is positive definite} and
is fully determined by the admissible tuples
\((a, b) \in \mathbb{R}^2\).

Using the simplified notation, the covariance matrix can be written as

\begin{equation}\protect\phantomsection\label{eq-sigma-ab}{
\Sigma =
\begin{pmatrix}
1 & a & b \\
a & 1 & b \\
b & b & 1
\end{pmatrix}.
}\end{equation}

The first leading principal minor always has a positive determinant
since \(\det(1) = 1 > 0\). The second leading principal minor is

\begin{equation}\protect\phantomsection\label{eq-second-principal-minor}{
\det
\begin{pmatrix}
1 & a \\
a & 1
\end{pmatrix}
=
1-a^2.
}\end{equation}

Therefore, the second condition requires

\begin{equation}\protect\phantomsection\label{eq-second-principal-minor-condition}{
1-a^2 > 0,
}\end{equation}

which is equivalent to

\begin{equation}\protect\phantomsection\label{eq-a-range}{
-1 < a < 1.
}\end{equation}

Finally, the determinant of the full matrix in
Equation~\ref{eq-sigma-ab} is

\begin{equation}\protect\phantomsection\label{eq-full-determinant-start}{
\det(\Sigma)
=
\begin{vmatrix}
1 & a & b \\
a & 1 & b \\
b & b & 1
\end{vmatrix}.
}\end{equation}

Subtracting the second row from the first row leaves the determinant
unchanged and yields

\begin{equation}\protect\phantomsection\label{eq-full-determinant-row-operation}{
\det(\Sigma)
=
\begin{vmatrix}
1-a & a-1 & 0 \\
a & 1 & b \\
b & b & 1
\end{vmatrix}.
}\end{equation}

Factoring out \((1-a)\) from the first row gives

\begin{equation}\protect\phantomsection\label{eq-full-determinant-factorized}{
\det(\Sigma)
=
(1-a)
\begin{vmatrix}
1 & -1 & 0 \\
a & 1 & b \\
b & b & 1
\end{vmatrix}.
}\end{equation}

The cofactor expansion results in

\begin{equation}\protect\phantomsection\label{eq-full-determinant-expansion}{
\det(\Sigma)
=
(1-a)
\left[
\begin{vmatrix}
1 & b \\
b & 1
\end{vmatrix}
+
\begin{vmatrix}
a & b \\
b & 1
\end{vmatrix}
\right].
}\end{equation}

Evaluating the two \(2 \times 2\) determinants yields

\begin{equation}\protect\phantomsection\label{eq-full-determinant-evaluated}{
\det(\Sigma)
=
(1-a)
\left[
(1-b^2) + (a-b^2)
\right].
}\end{equation}

After simplification,

\begin{equation}\protect\phantomsection\label{eq-full-determinant}{
\det(\Sigma)
=
(1-a)(1+a-2b^2).
}\end{equation}

By Equation~\ref{eq-a-range}, the factor \((1-a)\) is strictly positive.
Consequently, the third condition of Sylvester's criterion is

\begin{equation}\protect\phantomsection\label{eq-third-principal-minor-condition}{
1+a-2b^2 > 0.
}\end{equation}

Solving for \(b\) gives

\begin{equation}\protect\phantomsection\label{eq-b-condition-squared}{
b^2 < \frac{1+a}{2},
}\end{equation}

Combining Equation~\ref{eq-a-range} and
Equation~\ref{eq-b-condition-squared}, the set of positive definite
covariance matrices of the form Equation~\ref{eq-general-Sigma-3d} is
characterized by

\begin{equation}\protect\phantomsection\label{eq-positive-definite-region}{
\left\{
(a,b)\in\mathbb{R}^2 :
-1<a<1,\;
b^2<\frac{1+a}{2}
\right\}.
}\end{equation}

\subsubsection{On the Bayes Classifier with idiosyncratic
shocks}\label{on-the-bayes-classifier-with-idiosyncratic-shocks}

\paragraph{Decision
boundary}\label{sec-dec-boundary-perf-bayes-idyiosincratic-shocks}

The decision boundary between two classes is defined as the equal
probability region for both. For our case with two classes there is only
one boundary and it is given by \[
\mathbb{P}(Y=b | X=x) = \mathbb{P}(Y=g | X=x)
\] \[
\Leftrightarrow \mathbb{P}(Y=b | X=x) = \frac{1}{2}
\]

Plugging in Equation~\ref{eq-pd-with-shock-perf-bayes} we get \[
\pi_{\epsilon} \, p_{\tilde{Y} | X} (b, x) + (1 - \pi_{\epsilon}) \pi_{\varepsilon, b}
= \frac{1}{2}
\]
\begin{equation}\protect\phantomsection\label{eq-decision-boundary-in-terms-of-clean}{
\Leftrightarrow p_{\tilde{Y} | X} (b, x) = \underbrace{\frac{\frac{1}{2} - (1 - \pi_{\epsilon}) \pi_{\varepsilon, b}}{\pi_{\epsilon}}}_{=: \xi}
}\end{equation}

Note that Equation~\ref{eq-decision-boundary-in-terms-of-clean} shows
how considering the idiosyncratic shocks only moves the decision
boundary. If there were no shocks (\(\pi_\epsilon = 1\)) or their effect
was symmetric \(\left(\pi_{\varepsilon, b} = \frac{1}{2} \right)\) then
\(\xi = \frac{1}{2}\), which equals the decision boundary based on a
classifier ignoring the shocks.

Using Equation~\ref{eq-posterior-of-Y}, Equation~\ref{eq-marginal-of-X}
and Equation~\ref{eq-marginal-of-Y} we get \[
p_{\tilde{Y} | X} (b, x) = \frac{\phi_{X_b}(x) \pi_b}{\pi_b \phi_{X_b}(x) + \pi_g \phi_{X_g}(x)} = \xi
\] and using \(\pi_g = 1 - \pi_b\) \[
\Leftrightarrow  \pi_b(1 - \xi)\phi_{X_b}(x) = \xi (1 - \pi_b) \phi_{X_g}(x)
\] which allows for a odds ratio perspective
\begin{equation}\protect\phantomsection\label{eq-odds-ratio-shock}{
\Leftrightarrow \frac{\phi_{X_b}(x)}{\phi_{X_g}(x)} = \frac{\xi (1 - \pi_b)}{\pi_b(1 - \xi)} =: c
}\end{equation}

Since we assume equal covariances matrices we can get from taking loggs
the familiar result (see e.g. \textcite{hastie2009elements}, p.~108) \[
x^\top \Sigma^{-1} (\mu_b - \mu_g) - \underbrace{\frac{1}{2} (\mu_b + \mu_g)}_{=:\bar{\mu}} \Sigma^{-1}(\mu_b - \mu_g) =
\log c
\] After factorizing and using the symmetry of \(\Sigma^{-1}\) we get
\begin{equation}\protect\phantomsection\label{eq-decision-boundary-equal-covs}{
(\mu_b - \mu_g)^\top \Sigma^{-1} (x - \bar{\mu})  = \log c
}\end{equation}

\paragraph{Bayes error rate}\label{sec-derivation-bayes-rate}

A perfect Bayes classifier considering the idiosyncratic shock as
defined in Equation~\ref{eq-pd-with-shock-perf-bayes} has an analytical
solution for the Bayes error. To get to this expression, I use the
results of \textcite{Ripley_1996}.

For any \(\omega \in \Omega\), the decision to classify an applicant as
bad with characteristics \(x := X(\omega)\) occurs if \[
\mathbb{P}(Y=b | X = x) > \mathbb{P}(Y=g | X = x)
\]

which following the same steps as before we get an expression like
Equation~\ref{eq-decision-boundary-equal-covs} that we can use to also
define a random variable \[
A(\omega) := (\mu_b - \mu_g)^\top \Sigma^{-1} (x - \bar{\mu}) >
\log c
\] where, following \textcite{Ripley_1996}, \[
A | Y = b \sim \mathcal{N}\left(\frac{1}{2} \delta^2, \delta^2 \right), \,
A | Y = g \sim \mathcal{N}\left(-\frac{1}{2} \delta^2, \delta^2 \right)
\] with \[
\delta = \sqrt{(\mu_b - \mu_g)^\top \Sigma^{-1} (\mu_b - \mu_g)}
\]

Therefore the decision function \(d_B\) of the bayes classifier can be
written as \[
d_B(x) = \begin{cases}
    b & A > \log c \\
    g & A \leq \log c
\end{cases}
\]

The bayes error rate \(E_B\) is then given by
\begin{equation}\protect\phantomsection\label{eq-bayes-error-rate-formula}{
\begin{aligned}
E_B &= 
    \mathbb{P}(Y = b) \mathbb{P}(A \leq \log c | Y = b) +
    \mathbb{P}(Y = g)  \mathbb{P}(A > \log c | Y = g)
\\
&=
    \pi_b \Phi\left(-\frac{1}{2}\delta + \frac{1}{\delta} \log c\right) +
    (1-\pi_b) \Phi\left(-\frac{1}{2}\delta - \frac{1}{\delta} \log c\right)
\end{aligned}
}\end{equation}

\subsubsection{A possible route towards a formal justification of
surprise as a valid
signal}\label{a-possible-route-towards-a-formal-justification-of-surprise-as-a-valid-signal}

The argument in Section Section~\ref{sec-an-unsurprising-setting} is
intentionally heuristic. Nevertheless, it is useful to outline how one
could attempt to formalize the claim that surprise should vanish in a
stable MAR environment, which is the sole purpose of this appendix.

Assume the conditions from Section
Section~\ref{sec-an-unsurprising-setting} hold. Define the acceptance
region at round \(t\) as

\[
A_t
:=
\left\{
x \in \mathbb{R}^{f_m}
\;\middle|\;
\hat{s}_t(x) > c_t
\right\}.
\]

The corresponding distribution of accepted applicants is then

\[
\mathbb{P}_{Y,X \mid X \in A_t}(E)
=
\frac{
    \mathbb{P}\left((Y,X)\in E, X\in A_t\right)
}{
    \mathbb{P}\left(X\in A_t\right)
}
\]

for measurable \(E \in 2^{E_Y} \otimes \mathcal{B}(\mathbb{R}^f)\).

Suppose additionally that the estimation method is unbiased in the sense
of this thesis, namely that given enough representative data, the
estimated scorecard \(\hat{s}_t\) provides a good approximation of a
scoring proportional to \(\mathbb{P}(Y=b\mid X)\). Under suitable
regularity conditions and growing sample sizes, one could then expect
the sequence of scorecards \(\hat{s}_t\) to stabilize around some
limiting scorecard \(s^* \propto \mathbb{P}(Y=b\mid X)\).

Likewise, if the cutoff selection criterion remains unchanged across
rounds, one could expect the sequence of cutoffs \(c_t\) to stabilize
around a limiting value \(c\). This motivates the introduction of the
limiting acceptance region

\[
A
:=
\left\{
x \in \mathbb{R}^{f_m}
\;\middle|\;
s^*(x) > c
\right\}.
\]

Heuristically, stabilization of \(\hat{s}_t\) and \(c_t\) suggests that
the acceptance regions \(A_t\) become increasingly similar to \(A\). As
a consequence, one would expect the distributions of accepted applicants

\[
\mathbb{P}_{Y,X \mid X \in A_t}
\]

to become increasingly similar across rounds as well.

Although turning this intuition into a rigorous theorem would require
additional assumptions on the estimation method, the applicants
distribution and the acceptance regions, the central idea is simple: in
a stationary MAR environment the acceptance mechanism itself should
eventually stabilize. Once this occurs, both the training data and
future accepted applicants are generated from nearly the same
distribution.

Under such circumstances, one would expect the performance estimated
through cross-validation on historical accepts to be close to the
performance observed on future accepted applicants, implying only small
values of the surprise signal \(\eta^{(t)}\).

\subsection{Detailed material}\label{detailed-material}

In this part of the appendix I will add some tables and graphs that are
helpful to better understand some of the analysis done in the thesis.

\newpage{}

\begin{landscape}

\subsubsection{Literature Review table}\label{literature-review-table}

The table \textcite{BuckerEtAl2013MNARRI}

\end{landscape}

\begin{table}

\caption{\label{tbl-lit-review}}

\centering{

\begin{tabular}{llllllllll}
\toprule
Reference & RI Approach & Key identifying assumptions & Assumptions explicit? & Theoretical/heuristic justification? & Infers Rejects label for evaluation? & Type of result & Concludes RI can (reliably) work? & Main challenge & Dynamic perspective \\
\midrule
\cite{HandHanley1993CanRIWork} & Estimation of $p(x)$ for further usage in EM based methods & Distributional assumptions, sample of reject region or experts judgement & $\checkmark$ & $\checkmark$ & No & Theoretical & Only with full sample of reject region & Economic costs of sample &  \\
\cite{Anderson2019BayesianNetForRI} & Extrapolation to rejects through bayesian network, then build scorecard on augemented sample & Knowledge and existence of causal non-cyclic relationships between features & $\checkmark$ &  & No & Empirical & Moderately & Reliability of causality and dimensionality &  \\
\cite{LiEtAl2017SemisupervisedSVM} & Rejects guide the decision boundary through low density regions in the features space & Low density regions of features space aligned with known labels separation &  &  & Yes & Empirical & Yes, evaluation method is unreliable & High dimensional feature space problematic &  \\
\cite{BuckerEtAl2013MNARRI} & Estimation of unconditional means of features and probability of reject & Label missingness be independent across observations & $\checkmark$ & $\checkmark$ & No & Theoretical and simulation & Yes, in theoretical and simulation based results & Independence condition is illogical in dynamic context, numerical problems &  \\
\cite{Banasik01032010} & Augmentation via reweighting using accepts rejects model & Accept reject model estimable & $\checkmark$ & $\checkmark$ & No & Empirical & No, maybe with high rejection rates & Augmentation is not adequate for RI &  \\
\bottomrule
\end{tabular}

}

\end{table}%


\printbibliography



\end{document}
